% =====================================================================
%  ETSI GS QKD 018: orchestrazione SDN di reti QKD multi-dominio
%  ---------------------------------------------------------------------
%  COMPILAZIONE (XeLaTeX o LuaLaTeX, NON pdfLaTeX):
%      latexmk -xelatex main.tex
%  oppure, a mano:
%      xelatex main.tex ; biber main ; xelatex main.tex ; xelatex main.tex
%  Il font matematico e' Cambria Math (unicode-math, \setmathfont),
%  installato di default su Windows e con Microsoft Office. Se manca
%  (per esempio su Overleaf) si usa automaticamente Latin Modern Math.
%  Su Overleaf: Menu -> Compiler -> XeLaTeX.
%
%  Struttura del progetto:
%      main.tex           documento principale
%      preamble.tex       pacchetti, font, ambienti
%      abstract/          abstract
%      chapters/          un file per sezione
%      figures/           figure TikZ
%      bibliography/      references.bib
% =====================================================================
\documentclass[11pt,a4paper]{article}
% =====================================================================
%  Preambolo: pacchetti, font, ambienti e stili
% =====================================================================

% ---------- Font e lingua (richiede XeLaTeX o LuaLaTeX) ----------
\usepackage{fontspec}
\usepackage{amsmath,amssymb}
\usepackage{amsthm}
\usepackage{unicode-math}
\usepackage[italian]{babel}
\usepackage{csquotes}
\addto\captionsitalian{\renewcommand{\abstractname}{Abstract}}

% Cambria Math per gli ambienti matematici; fallback se non installato
\IfFontExistsTF{Cambria Math}
  {\setmathfont{Cambria Math}}
  {\setmathfont{Latin Modern Math}}
% Per avere anche il testo in Cambria, decommentare:
% \IfFontExistsTF{Cambria}{\setmainfont{Cambria}}{}

% ---------- Impaginazione ----------
\usepackage[a4paper,margin=2.5cm]{geometry}
\usepackage{microtype}
\usepackage{booktabs,tabularx,array,longtable}
\usepackage{enumitem}
\usepackage[font=small,labelfont=bf]{caption}
\usepackage{xcolor}
\usepackage{listings}
\usepackage{tikz}
\usetikzlibrary{arrows.meta,positioning,fit,calc,shapes.geometric,backgrounds}
\usepackage{authblk}

% ---------- Bibliografia ----------
\usepackage[backend=biber,style=numeric-comp,sorting=none,maxbibnames=6]{biblatex}
\addbibresource{bibliography/references.bib}

\usepackage[hidelinks]{hyperref}
\usepackage{xurl}

% ogni sezione inizia su una pagina nuova
\usepackage{etoolbox}
\pretocmd{\section}{\clearpage}{}{}

\setlist{itemsep=2pt,topsep=4pt}
\setlength{\parskip}{3pt}
\setlength{\emergencystretch}{2em}
\newcolumntype{L}{>{\raggedright\arraybackslash}X}
% a capo consentito dopo il trattino basso negli identificatori \texttt
\AtBeginDocument{\renewcommand{\_}{\textunderscore\allowbreak}}

% ---------- Ambiente "definition" ----------
\theoremstyle{definition}
\newtheorem{definition}{Definizione}[section]

% ---------- Listati (YANG, alberi pyang, JSON, HTTP) ----------
\definecolor{codebg}{gray}{0.96}
\definecolor{codekw}{RGB}{0,70,140}
\definecolor{codecm}{RGB}{110,110,110}
\definecolor{codestr}{RGB}{150,40,40}
\lstdefinelanguage{yang}{
  morekeywords={module,namespace,prefix,import,organization,contact,
    description,revision,reference,yang-version,identity,base,typedef,type,
    identityref,grouping,uses,container,list,key,leaf,leaf-list,config,
    augment,rpc,input,output,notification,leafref,path,default,units,
    must,when,enum,enumeration,mandatory},
  sensitive=true,
  morecomment=[l]{//},
  morecomment=[s]{/*}{*/},
  morestring=[b]"
}
\lstdefinestyle{code}{
  basicstyle=\ttfamily\footnotesize,
  keywordstyle=\color{codekw}\bfseries,
  commentstyle=\color{codecm}\itshape,
  stringstyle=\color{codestr},
  backgroundcolor=\color{codebg},
  frame=single, framerule=0pt,
  columns=fullflexible, keepspaces=true,
  xleftmargin=4pt, xrightmargin=4pt,
  breaklines=true, upquote=true,
  aboveskip=6pt, belowskip=6pt
}
\lstset{style=code}
\renewcommand{\lstlistingname}{Listato}

% ---------- Stili TikZ comuni ----------
\tikzset{
  >=Stealth,
  box/.style={draw,rounded corners,minimum height=8mm,align=center,
              font=\small},
  lbl/.style={font=\scriptsize,midway,fill=white,inner sep=1pt},
  actor/.style={draw,rounded corners,fill=gray!8,minimum width=2.4cm,
                minimum height=7mm,align=center,font=\small},
  msg/.style={->,thick},
  ret/.style={->,dashed}
}


\title{\textbf{Orchestrazione SDN di reti QKD multi-dominio:\\
       l'interfaccia ETSI GS QKD 018}}
\author{Diego Carro}
\author{Bucci Felice Elia}
\affil{Corso di Cybersecurity, Politecnico di Torino\\Quantum Engineering}
\date{Novembre 2026}

\begin{document}
\maketitle

% =====================================================================
%  Abstract
% =====================================================================
\begin{abstract}
\noindent
La distribuzione quantistica di chiavi (QKD) permette a due parti di
concordare una chiave segreta con una sicurezza fondata sulle leggi della
fisica anziché sulla difficoltà computazionale. Portarla dal laboratorio
alla rete di un operatore richiede però di gestirla come qualunque altra
risorsa di rete: scoprirne la topologia, monitorarne lo stato, allocarne la
capacità ai servizi e coordinarla con le reti classiche. Il gruppo ETSI ISG
QKD ha affrontato il problema per gradi. La GS~QKD~002 individua i casi
d'uso e i requisiti di integrazione; la GS~QKD~015 definisce il modello con
cui un controller SDN vede e configura i nodi QKD di un dominio; la
GS~QKD~018 definisce l'interfaccia tra un \emph{SDN orchestrator} e il
controller di un dominio QKD, abilitando il provisioning end-to-end tra più
domini e tra rete QKD e rete di trasporto ottico. Questo lavoro ripercorre
la catena in modo ragionato: descrive le componenti di una rete SD-QKD e la
struttura di un modello YANG, analizza in dettaglio la GS~QKD~018 (casi
d'uso, modello informativo, moduli di topologia, inventario e connettività,
workflow multi-dominio, sicurezza), ne mostra l'analogia con il framework
ACTN dell'IETF (RFC~8453) e la confronta con tre esperienze reali: il
testbed multi-vendor KOREN di SK Telecom, un prototipo gerarchico basato su
ONOS sviluppato al Politecnico di Milano e la rete MadQCI di Madrid. Si
conclude con i limiti attuali dello standard e le direzioni di ricerca
aperte.

\medskip
\noindent\textbf{Parole chiave:} QKD, SDN, orchestrazione, ETSI GS QKD 015,
ETSI GS QKD 018, YANG, RESTCONF, multi-dominio, gestione delle chiavi.
\end{abstract}


\tableofcontents
\bigskip

% =====================================================================
\section{Stato dell'arte}
\label{sec:sota}
% =====================================================================

\paragraph{Fondamenti di QKD.}
La letteratura sulla QKD è ormai consolidata. La review di Pirandola
et al.~\cite{pirandola2020} copre i protocolli a variabili discrete (BB84 e
le sue varianti con \emph{decoy state}), quelli a variabili continue e le
dimostrazioni di sicurezza. Da qui emergono i due vincoli che guidano tutta
l'ingegneria di rete: il tasso di chiave decresce esponenzialmente con la
distanza e un singolo collegamento in fibra commerciale non supera le poche
centinaia di chilometri.

\paragraph{Dal collegamento alla rete.}
Mehic et al.~\cite{mehic2020} e Cao et al.~\cite{cao2022} descrivono il
passaggio dal collegamento punto-punto alla \emph{rete} QKD: architettura a
strati, \emph{trusted relay}, gestione e instradamento delle chiavi. Il
white paper GÉANT~\cite{geant2024} offre una sintesi leggibile dello stesso
quadro e degli standard ETSI e ITU-T. Il problema dell'instradamento, che
diventa centrale quando si orchestra una rete, è trattato
in~\cite{routing2024}.

\paragraph{Standardizzazione.}
Due organismi hanno prodotto specifiche complementari. L'ETSI ISG QKD ha
iniziato nel 2010 con i casi d'uso~\cite{qkd002} e ha poi definito le
interfacce verso le applicazioni~\cite{qkd004,qkd014}, verso il controller
SDN~\cite{qkd015}, verso l'orchestrator~\cite{qkd018} e tra sistemi di
gestione delle chiavi~\cite{qkd020}. L'ITU-T, con la serie Y.3800, ha
definito l'architettura di riferimento delle reti QKD
(QKDN)~\cite{ituy3800}, il loro controllo e la loro
gestione~\cite{ituy3804} e il controllo tramite SDN~\cite{ituy3805}.

\paragraph{SDN applicato alla QKD.}
Il paradigma \emph{Software-Defined Networking}~\cite{kreutz2015} separa il
piano di controllo dal piano dati e accentra le decisioni in un controller
programmabile. Il gruppo UPM/Telefónica ha mostrato per primo
l'integrazione di risorse QKD in una rete ottica controllata via SDN, sia
sperimentalmente~\cite{aguado2016} sia in un tutorial
sistematico~\cite{aguado2020}; da questo filone nasce la GS~QKD~015.

\paragraph{Orchestrazione e multi-dominio.}
L'interfaccia orchestrator–controller è specificata dalla
GS~QKD~018~\cite{qkd018} e dai relativi moduli YANG~\cite{qkd018yang}. Il
primo caso reale documentato è quello di Sim et al.~\cite{sim2023}
(SK~Telecom), che orchestrano due domini di vendor diversi sul testbed
KOREN. La tesi di Vicario~\cite{vicario2025} realizza un prototipo
gerarchico completo con ONOS e RESTCONF. La rete MadQCI~\cite{madqci2024}
usa estesamente la GS~QKD~015 ma indica l'orchestrazione GS~QKD~018 come
attività ancora in corso. Lavori più recenti propongono applicazioni SDN per
la riconfigurazione dinamica dei link~\cite{mendez2026}, criteri di
selezione dei cammini e instradamento inter-dominio che preserva la
riservatezza della topologia~\cite{laschet2026} e modelli ibridi di
consegna delle chiavi tra reti regionali~\cite{interconnect2026}.

\paragraph{Contributo.}
Rispetto a queste fonti, il presente lavoro offre una lettura unitaria e
ragionata della catena GS~QKD~002 $\to$ 015 $\to$ 018: mostra come i
requisiti dei casi d'uso si traducano nelle componenti di una rete SD-QKD,
spiega la struttura dei modelli YANG che le descrivono e analizza come la
GS~QKD~018 li astragga per l'orchestrazione multi-dominio.

% =====================================================================
\section{Introduzione}
\label{sec:intro}
% =====================================================================

Questa sezione costruisce il vocabolario necessario per leggere le
specifiche ETSI: parte dalla QKD come primitiva crittografica, passa alla
rete QKD con i suoi nodi fidati e la sua architettura a strati, richiama i
concetti essenziali di SDN e arriva infine ai due documenti che precedono la
GS~QKD~018: la GS~QKD~002 (§\ref{sec:qkd002}) e la GS~QKD~015
(§\ref{sec:qkd015}).

\paragraph{La distribuzione quantistica di chiavi.}

\begin{definition}[QKD]
La \emph{Quantum Key Distribution} è una procedura che, tramite il
trasporto di stati quantistici, consente a due parti remote di concordare
stringhe di bit segrete condivise, con una sicurezza basata
sull'entanglement o sull'impossibilità di clonare o misurare perfettamente
stati quantistici sconosciuti~\cite{qkd018}.
\end{definition}

L'idea si capisce bene con BB84. Alice codifica bit casuali su singoli
fotoni in una di due basi scelte a caso; Bob misura in una base anch'essa
casuale; su un canale classico \emph{autenticato} i due confrontano le basi
e tengono solo i bit misurati nella base giusta (\emph{sifting}). Un
intercettatore che misura i fotoni introduce inevitabilmente errori,
quantificati dal \emph{Quantum Bit Error Rate}
\begin{equation}
  \mathrm{QBER} = \frac{N_{\text{errori}}}{N_{\text{sifted}}} .
\end{equation}
Se il QBER resta sotto una soglia (circa l'$11\%$ per BB84 contro attacchi
collettivi), correzione d'errore e \emph{privacy amplification} distillano
una chiave segreta; altrimenti il blocco viene scartato. La grandezza che
interessa alla rete è il \emph{secret key rate} (SKR), cioè i bit di chiave
sicura prodotti al secondo. Poiché la trasmittanza della fibra decade come
$\eta = 10^{-\alpha L/10}$, con $\alpha \approx 0{,}2$~dB/km, anche l'SKR
cala rapidamente con la lunghezza $L$ del collegamento.

Due osservazioni, che ritroveremo nei casi d'uso della GS~QKD~002, sono
fondamentali. La QKD \emph{distribuisce} chiavi ma non cifra né autentica:
va sempre combinata con un cifrario (one-time pad o AES) e con un
meccanismo di autenticazione. Inoltre il canale classico deve essere
autenticato fin dal primo scambio, il che richiede un piccolo segreto
pre-condiviso: la QKD è quindi, a rigore, un meccanismo di \emph{key
growing}.

\paragraph{Dal link alla rete: il trusted node.}
Un link QKD copre al massimo qualche centinaio di chilometri e non esistono
ancora ripetitori quantistici pratici. Le reti reali collegano quindi i
link in serie tramite nodi fidati.

\begin{definition}[Trusted node e key relay]
Un \emph{trusted node} è un nodo, collocato in un perimetro fisicamente
protetto, che condivide chiavi QKD con entrambi i vicini. Per trasportare
una chiave $K$ da $A$ a $C$ passando per $B$, il nodo $A$ invia
$c = K \oplus K_{AB}$; $B$ calcola $K = c \oplus K_{AB}$ e inoltra
$c' = K \oplus K_{BC}$. Ogni hop è cifrato con \emph{one-time pad}, ma $K$
compare in chiaro dentro ciascun nodo intermedio, che deve quindi essere
fidato.
\end{definition}

Si noti che il relay consuma chiave: per trasportare $|K|$ bit lungo un
cammino $p$ ogni link $\ell\in p$ spende $|K|$ bit della propria chiave
locale. Questo semplice fatto rende la chiave una \emph{risorsa di rete} da
pianificare, ed è il motivo per cui le specifiche successive parlano di
tassi disponibili, consumi e cammini.

\paragraph{L'architettura a strati ITU-T.}
La raccomandazione ITU-T Y.3800~\cite{ituy3800} organizza una rete QKD
(QKDN) in strati sovrapposti (Fig.~\ref{fig:itu}): il \emph{quantum layer}
con i moduli QKD; il \emph{key management layer} con i key manager che
memorizzano e inoltrano le chiavi; il \emph{control layer}, che prende
decisioni operative (instradamento, relay, configurazione dei link); il
\emph{management layer}, che supervisiona il tutto; e il \emph{service
layer}, dove le applicazioni consumano le chiavi. In questo schema si
distinguono due tipi di collegamento.

\begin{definition}[QKD link e KM link]
Un \emph{QKD link} connette due moduli QKD ed è composto da canale
quantistico e canale classico: è il luogo in cui la chiave viene
\emph{generata}. Un \emph{KM link} connette due \emph{key manager} ed è il
canale classico, protetto con chiavi QKD, lungo il quale la chiave viene
\emph{inoltrata} (key relay).
\end{definition}

Allo stesso modo si distinguono due ruoli spesso confusi. Il \emph{QKDN
controller} agisce sulla rete in tempo reale: decide i cammini di relay,
attiva o disattiva link, reagisce ai guasti. Il \emph{QKDN manager} ha un
compito di supervisione: raccoglie informazioni FCAPS (\emph{Fault,
Configuration, Accounting, Performance, Security}) e applica politiche. Le
specifiche ETSI che vedremo si collocano nel piano di controllo: la
GS~QKD~015 tra controller e nodi, la GS~QKD~018 tra controller e
orchestrator.

% Figura: architettura a strati di una QKDN (ITU-T Y.3800)
\begin{figure}[t]
\centering
\begin{tikzpicture}[
  layer/.style={draw,rounded corners,minimum width=11cm,minimum height=7mm,
                align=center,font=\small},
  side/.style={font=\scriptsize,align=left,anchor=west}]
  \node[layer,fill=gray!10]   (srv) at (0,0)    {Service layer: applicazioni / SAE};
  \node[layer,fill=red!8]     (mgm) at (0,-1.0) {Management layer: QKDN manager (FCAPS)};
  \node[layer,fill=blue!8]    (ctl) at (0,-2.0) {Control layer: QKDN controller (SDN)};
  \node[layer,fill=green!10]  (km)  at (0,-3.0) {Key management layer: key manager (KME), KM link};
  \node[layer,fill=yellow!15] (q)   at (0,-4.0) {Quantum layer: moduli QKD, QKD link};
  \node[side] at (5.7,0)    {QKD 004 / 014};
  \node[side] at (5.7,-1.5) {QKD 018\\(orchestrazione)};
  \node[side] at (5.7,-2.5) {QKD 015\\(controllo)};
\end{tikzpicture}
\caption[Architettura a strati di una rete QKD (ITU-T Y.3800)]{Architettura a strati di una rete QKD secondo ITU-T
Y.3800~\cite{ituy3800}, con a destra la collocazione indicativa delle
specifiche ETSI discusse nel testo.}
\label{fig:itu}
\end{figure}


Le applicazioni non vedono la parte quantistica. Un \emph{Secure
Application Entity} (SAE) chiede chiavi al \emph{Key Management Entity}
(KME) del proprio nodo tramite le API GS~QKD~004~\cite{qkd004} o
GS~QKD~014~\cite{qkd014}; due SAE che vogliono comunicare ricevono la
stessa chiave dai rispettivi KME, identificata da un \emph{key ID}.

\paragraph{Software-Defined Networking.}
In una rete tradizionale ogni apparato decide da solo come instradare il
traffico. SDN~\cite{kreutz2015} separa invece tre piani: il \emph{piano
dati}, fatto di dispositivi che eseguono; il \emph{piano di controllo},
accentrato in un controller con la vista globale della rete; il
\emph{piano di gestione e applicazione}, che esprime gli obiettivi. Il
controller parla con i dispositivi tramite un'interfaccia
\emph{southbound} (SBI) e offre i propri servizi verso l'alto tramite
un'interfaccia \emph{northbound} (NBI). Quando i domini sono più di uno, si
aggiunge sopra i controller un \emph{orchestrator}, che coordina servizi
end-to-end senza conoscere i dettagli di ciascun dominio.

Una rete QKD ha tutte le caratteristiche che rendono utile SDN: risorse
scarse e variabili (la chiave), vincoli fisici stringenti (attenuazione,
coesistenza con canali classici), apparati di vendor diversi e domande che
cambiano nel tempo. Le due specifiche che seguono descrivono perché e come
l'ETSI ha applicato questo paradigma alla QKD.

% =====================================================================
\subsection{ETSI GS QKD 002: i casi d'uso}
\label{sec:qkd002}
% =====================================================================

La GS~QKD~002 V1.1.1 (giugno 2010)~\cite{qkd002} è il primo documento
pubblicato dall'ISG QKD. Il suo obiettivo non è tecnico in senso stretto:
prima di standardizzare \emph{come} realizzare la QKD, il gruppo voleva
chiarire \emph{dove} serve, \emph{chi} la usa e \emph{quali requisiti}
deve rispettare per entrare nelle infrastrutture esistenti. Il documento
ragiona in tre passi: colloca la QKD tra le primitive crittografiche,
individua i livelli dello stack di rete in cui può essere inserita e
descrive sei casi d'uso con i relativi requisiti. Seguire questo
ragionamento è utile perché mostra da dove nascono le esigenze che, dodici
anni dopo, porteranno alle GS~QKD~015 e~018.

\subsubsection{La QKD come primitiva crittografica}

Il documento parte da un sistema QKD tipico: due unità, di solito montate
in rack da 19", collegate da un canale quantistico in fibra o in spazio
libero lungo fino a circa $100$~km, che producono alle due estremità la
stessa sequenza di bit casuali. La clausola 4 chiarisce subito che questo
risolve \emph{un solo} problema crittografico.

Un sistema sicuro ha bisogno di tre primitive: \emph{cifratura},
\emph{distribuzione delle chiavi} e \emph{autenticazione dei messaggi}. La
QKD copre solo la seconda; per la cifratura servono one-time pad o cifrari
simmetrici come AES, per l'autenticazione servono MAC (per esempio basati
su hashing universale di tipo Wegman–Carter) o firme digitali. Ne segue che
la QKD è un \emph{complemento} alla crittografia classica, non un suo
sostituto, e che la sua sicurezza complessiva dipende anche dalle altre
componenti del sistema. Il documento sottolinea inoltre che il canale
pubblico della QKD va autenticato con un segreto pre-distribuito.

\subsubsection{Dove si inserisce la QKD nello stack di rete}

La clausola 6 descrive in quale livello dello stack una chiave QKD può
essere consumata. La scelta del livello determina chi deve essere
modificato e quanto l'integrazione è trasparente.

\begin{description}[style=nextline,leftmargin=1.2em]
  \item[Livello di collegamento (L2)] La chiave QKD alimenta lo scambio di
        chiavi di PPP, le chiavi di IEEE~802.1AE (MACsec) oppure un
        \emph{QKD link encryptor}, cioè un cifratore di linea che consuma
        direttamente le chiavi prodotte dal link e crea un tunnel cifrato
        tra due siti. È l'opzione più trasparente per le applicazioni.
  \item[Livello di rete (L3)] Una versione modificata di IKE usa la chiave
        QKD come segreto condiviso per IPsec.
  \item[Livello di trasporto (L4)] La chiave QKD sostituisce le chiavi di
        sessione di TLS/SSL o viene usata al posto degli HMAC per
        l'autenticazione.
  \item[Livello applicativo] Le applicazioni usano la chiave per
        autenticare gli utenti o cifrare il contenuto scambiato tra le
        proprie istanze.
\end{description}

Il punto concettuale è che, qualunque sia il livello, serve un componente
che \emph{consuma} chiavi e uno che le \emph{fornisce}: è la separazione
che più tardi diventerà SAE e KME, con un'interfaccia standard
(GS~QKD~004 e~014) tra i due.

\subsubsection{I sei casi d'uso}

La clausola 7 descrive sei casi d'uso con una struttura fissa: obiettivo,
descrizione, attori, benefici, considerazioni operative e di qualità del
servizio. Li riassumiamo seguendo un filo logico: si passa dal
collegamento singolo alla rete, e dalla rete privata alla rete di un
operatore.

\paragraph{UC1 – Offsite backup / business continuity.}
Un'impresa replica i dati da un sito primario a un sito di backup
(backup giornalieri, snapshot, mirroring di SAN, replica di database). Un
singolo link QKD punto-punto alimenta un cifratore di linea su Ethernet o
Fibre Channel; la fibra può essere dedicata (\emph{dark fibre}) o condivisa
in WDM o TDM. I requisiti emersi sono: rinnovo della chiave tanto frequente
quanto il tasso di generazione consente, gestione delle fasi di
ricalibrazione durante trasferimenti continui, un segreto iniziale
all'installazione e alta velocità di cifratura. Il documento nota che la
sicurezza complessiva non è incondizionata, perché dipende anche dal
cifratore.

\paragraph{UC2 – Rete metropolitana aziendale.}
Più sedi di un'impresa o di un ente pubblico sono collegate ai data
center, attraverso i quali passa anche il traffico tra sedi. Ogni
collegamento è protetto da un QKD link encryptor che cifra e autentica a
livello~2, in modo trasparente per le applicazioni. Emergono due esigenze
nuove: link diversi hanno volumi e tempi di traffico diversi, quindi
\emph{requisiti di QoS diversi}, e i trasferimenti possono essere
pianificati, periodici o attivati da soglie. La variante
\emph{QKD Secured Key Server} prevede un server centrale che distribuisce
chiavi e certificati ai client dell'area metropolitana su canali cifrati
con one-time pad e autenticati con Wegman–Carter: è un primo esempio di
servizio di chiavi offerto da una rete.

\paragraph{UC3 – Controllo di infrastrutture critiche.}
Il caso protegge il traffico SCADA di un'infrastruttura critica, con
l'esempio di una rete ferroviaria: qui conta più l'\emph{autenticità} dei
comandi (scambi, segnali) che la loro riservatezza. I nodi della WAN privata
stanno in stazioni con perimetro controllato, mentre la fibra è fuori dal
perimetro. Il documento propone due architetture: link QKD singoli tra
nodi, con buffer e inoltro delle chiavi gestiti nei nodi, oppure una
\emph{rete QKD con gestione delle chiavi integrata} e un'interfaccia
dedicata verso le applicazioni. Si chiedono la massima disponibilità, con
precauzioni per ricalibrazioni, guasti e denial of service, e la
possibilità di inizializzare nuovi link ridondanti con segreti ottenuti
lungo un cammino alternativo. È il primo caso in cui compaiono, insieme,
key relay, gestione centralizzata e ridondanza dei cammini.

\paragraph{UC4 – Protezione del backbone.}
Una fibra del backbone metropolitano di un operatore, tipicamente un anello
WDM con OADM o ROADM, trasporta un canale quantistico accanto ai canali
dati. Le chiavi proteggono il traffico di controllo e gestione o vengono
offerte ai clienti. Gli attori sono due: chi gestisce l'infrastruttura e
chi offre servizi sopra di essa; i nodi del backbone sono fidati. I
requisiti sono i più vicini a quelli di un operatore: il sistema QKD non
deve ridurre il throughput del sistema WDM di oltre il $10\%$, deve essere
compatibile con fibre, giunti e connettori comuni, la sua gestione deve
\emph{integrarsi in quella del sistema WDM} e installazione e manutenzione
devono essere possibili a personale senza competenze di fisica quantistica.

\paragraph{UC5 – Rete di accesso ad alta sicurezza.}
In una rete ottica passiva (PON) per Fiber-to-the-Home, il traffico verso
gli utenti è trasmesso in broadcast a tutte le ONU, quindi va cifrato per
impedire che un utente legga i dati di un altro. Un OLT serve da $32$ a
$128$ ONU e il canale quantistico segue lo stesso percorso di quello
classico. La configurazione è fortemente asimmetrica: un solo rivelatore
all'OLT, più costoso, serve molte sorgenti alle ONU. Le distanze brevi
permettono tassi elevati, necessari perché la capacità di generazione è
condivisa tra tutti gli utenti. La variante \emph{QKD Authenticated Sensor
Network} applica lo stesso schema a una stazione centrale con molti
sensori e attuatori.

\paragraph{UC6 – Servizi a lunga distanza (long-haul).}
Due stazioni a terra $A$ e $B$ molto distanti, per esempio sulle due
sponde di un oceano, vogliono condividere una chiave senza fidarsi di nodi
intermedi a terra. Un nodo mobile $C$ (satellite, aereo o piattaforma ad
alta quota) stabilisce un segreto $a$ con $A$ e poi un segreto $b$ con $B$,
e invia a $B$ il valore $a\oplus b$ su canale classico; $B$ ricava $a$. È
un key relay in cui il nodo fidato è $C$. I satelliti in orbita bassa
operano a circa $300$–$800$~km di quota e passano sopra le stazioni circa
una volta al giorno; l'orbita geostazionaria non è praticabile. Meteo e
atmosfera possono ridurre o annullare il tasso, quindi servono canali
ridondanti.

La Tab.~\ref{tab:qkd002} riassume i sei casi.

\begin{table}[h]
\centering\small
\caption{I casi d'uso della GS~QKD~002 e gli elementi che anticipano le
specifiche successive.}
\label{tab:qkd002}
\begin{tabularx}{\textwidth}{@{}l L L@{}}
\toprule
\textbf{Caso} & \textbf{Configurazione} & \textbf{Esigenza che anticipa} \\
\midrule
UC1 Backup & link punto-punto + cifratore di linea & interfaccia chiave–cifratore \\
UC2 Metro aziendale & più link, server di chiavi & QoS per applicazione \\
UC3 Infrastrutture critiche & rete con key relay e gestione integrata & KMS, ridondanza, controllo centralizzato \\
UC4 Backbone & canale quantistico in WDM & coesistenza con la rete ottica, gestione integrata \\
UC5 Accesso PON & punto-multipunto & condivisione della capacità di chiave \\
UC6 Long-haul & relay tramite nodo mobile & trusted node, ridondanza \\
\bottomrule
\end{tabularx}
\end{table}

\subsubsection{Che cosa lascia in eredità}

La GS~QKD~002 è un documento di prima edizione: la clausola sui requisiti
generali e le caratteristiche funzionali dei singoli casi sono indicate
come non ancora disponibili. Il suo valore sta nell'aver fissato le
domande. Letti in sequenza, i casi d'uso mostrano che la QKD utile non è il
singolo link ma una \emph{rete}, che richiede: un sistema di gestione delle
chiavi che le memorizzi e le inoltri (UC3, UC6); un modo per esprimere
requisiti di servizio diversi per applicazioni diverse (UC2); la convivenza
con la rete ottica dell'operatore, con una gestione integrata in quella
esistente e utilizzabile da personale non specializzato (UC4). Le prime due
esigenze troveranno risposta nella GS~QKD~015, la terza, nella forma del
coordinamento tra rete QKD e rete di trasporto, nella GS~QKD~018.

% =====================================================================
\subsection{ETSI GS QKD 015: il controller SDN e la rete SD-QKD}
\label{sec:qkd015}
% =====================================================================

La GS~QKD~015 (V1.1.1 nel 2021, V2.1.1 nell'aprile 2022)~\cite{qkd015}
definisce l'interfaccia di controllo tra un controller SDN e i nodi di una
rete QKD, cioè la \emph{southbound interface} del piano di controllo QKD.
La seconda edizione, quella considerata qui, è stata allineata alla
GS~QKD~018: i moduli YANG della 018 importano direttamente quelli della 015.
Capire bene la 015 è quindi il prerequisito per capire la 018.

La sezione procede in quattro passi: perché serve un controllo software;
quali sono, una per una, le componenti di una rete SD-QKD; come è fatto in
generale un modello YANG; come la 015 usa YANG per descrivere il nodo, con
i workflow e le considerazioni di sicurezza che ne derivano.

% ---------------------------------------------------------------------
\subsubsection{Perché un controllo software-defined}

La specifica parte da una constatazione: integrare la QKD in una rete di
telecomunicazioni è oggi un'operazione manuale. Ogni nuovo sistema QKD
viene installato a mano, di solito su una fibra scura dedicata che collega
due siti, formando un link fisso. Questo approccio ha tre limiti. Non
scala, perché ogni modifica richiede intervento umano. Spreca risorse,
perché la capacità di chiave di un link non può essere spostata dove serve.
Infine non dialoga con la rete ottica, che invece è sempre più spesso
disaggregata e controllata da software.

L'architettura proposta permette a un controller di orchestrare
centralmente le risorse QKD: ottimizzare l'allocazione di chiave per link
in base alla domanda, creare automaticamente collegamenti diretti o
multi-hop, coordinare canali quantistici e classici anche quando
condividono la stessa fibra. In particolare l'astrazione definita dalla
specifica mira a:
\begin{itemize}
  \item gestire centralmente le risorse QKD in base a domande, capacità e
        stato della rete;
  \item aggregare i sistemi QKD di uno stesso perimetro sotto un'unica
        gestione delle chiavi, per riconoscere meglio le domande;
  \item nascondere la complessità dei parametri di basso livello di ogni
        tecnologia, delegandoli al nodo;
  \item ottimizzare la distribuzione dei link in base alla QoS delle
        applicazioni e alle statistiche di generazione;
  \item coordinare canali quantistici e classici, condividano o meno lo
        stesso mezzo fisico.
\end{itemize}
Sono, in forma tecnica, le esigenze emerse dai casi UC2, UC3 e UC4 della
GS~QKD~002.

% ---------------------------------------------------------------------
\subsubsection{Le componenti di una rete SD-QKD}

Conviene costruire la rete dal basso, aggiungendo un'astrazione alla volta.
La Fig.~\ref{fig:sdqkd} mostra il risultato.

\paragraph{1. Modulo, sistema e link QKD.}
L'unità fisica di base è il \emph{modulo QKD}: l'insieme di hardware e
software che esegue una parte di un protocollo QKD. Due moduli collegati
formano un \emph{sistema QKD}, e il collegamento tra loro (canale
quantistico più canale classico) è il \emph{QKD link}. Fin qui siamo nel
quantum layer della Fig.~\ref{fig:itu}: i moduli sono specifici del vendor
e della tecnologia (variabili discrete o continue, codifica, lunghezza
d'onda), e non sono interoperabili tra loro.

\paragraph{2. Interfaccia QKD.}
Il controller non ha bisogno di sapere come funziona internamente un
modulo, ma solo che cosa offre alla rete. La 015 introduce quindi
un'astrazione.

\begin{definition}[Interfaccia QKD]
Un'\emph{interfaccia QKD} (\texttt{qkdi}) è la rappresentazione astratta di
un modulo QKD all'interno di un nodo: ne espone solo gli attributi rilevanti
per la rete, come tecnologia, modello, intervallo di lunghezze d'onda,
attenuazione massima supportata e punto di attacco a uno switch ottico.
\end{definition}

L'analogia è quella delle porte di uno switch: un nodo con tre moduli QKD
appare al controller come un apparato con tre interfacce. Poiché moduli di
vendor diversi non interoperano, il modello impone di dichiarare tecnologia
e modello di ciascuna interfaccia: il controller abbinerà solo interfacce
compatibili.

\paragraph{3. Nodo SD-QKD.}
Le interfacce che stanno nello stesso luogo sicuro vengono raggruppate.

\begin{definition}[Nodo SD-QKD]
Un \emph{nodo SD-QKD} è l'aggregazione di uno o più moduli QKD situati in
un luogo sicuro e gestiti da un unico \emph{SDN agent}, che li espone al
controller come un singolo elemento di rete con più interfacce.
\end{definition}

La specifica fissa una serie di requisiti che danno al nodo la sua forma.
Un nodo SD-QKD \emph{deve}: comprendere almeno un modulo QKD, esposto come
interfaccia; trovarsi in un luogo sicuro (un luogo può ospitare più nodi);
contenere un \emph{sistema di gestione delle chiavi} che aggrega il
materiale prodotto da tutte le associazioni di chiave; offrire alle
applicazioni un \emph{unico punto di accesso} alle chiavi tramite API;
collegarsi ad almeno un controller SDN con protocolli standard, per
discovery, controllo e telemetria. Inoltre \emph{dovrebbe} esporre al
controller le informazioni su applicazioni, interfacce, link e requisiti
del canale classico (per esempio l'attenuazione accettabile).

All'interno del nodo si riconoscono quindi quattro componenti logiche:
\begin{itemize}
  \item i \textbf{moduli QKD}, che generano chiave con i nodi remoti;
  \item il \textbf{key manager} (KMS), che riceve le chiavi dai moduli,
        le memorizza in uno o più \emph{key store} logici, le inoltra in
        caso di relay e registra le applicazioni con la loro QoS;
  \item l'\textbf{interfaccia verso le applicazioni}, cioè la key
        delivery API (GS~QKD~004 o~014), raggiungibile a un indirizzo
        IP e porta noti;
  \item l'\textbf{SDN agent}, che traduce lo stato di tutto ciò in un
        modello YANG, lo espone al controller, riceve configurazioni e
        invia notifiche.
\end{itemize}

% Figura: rete SD-QKD e struttura interna di un nodo (adattata da GS QKD 015, Fig. 1)
\begin{figure}[t]
\centering
\begin{tikzpicture}[
  comp/.style={draw,rounded corners,align=center,font=\scriptsize,
               minimum height=6mm},
  small/.style={draw,circle,fill=green!12,minimum size=9mm,font=\scriptsize,
                align=center}]
  % controller
  \node[box,fill=blue!8,minimum width=4cm] (ctl) at (-0.6,4.3) {SDN controller (SD-QNC)};
  % nodo dettagliato
  \node[comp,fill=yellow!20,minimum width=1.4cm] (m1) at (-2.2,0) {modulo\\QKD 1};
  \node[comp,fill=yellow!20,minimum width=1.4cm] (m2) at (-0.6,0) {modulo\\QKD 2};
  \node[comp,fill=yellow!20,minimum width=1.4cm] (m3) at (1.0,0) {modulo\\QKD 3};
  \node[comp,fill=green!12,minimum width=4.6cm] (kms) at (-0.6,1.0) {key manager (KMS, key store)};
  \node[comp,fill=blue!6,minimum width=1.9cm] (agent) at (-1.95,2.0) {SDN agent};
  \node[comp,fill=gray!10,minimum width=2.4cm] (app) at (0.75,2.0) {applicazioni\\(QKD 004/014)};
  \begin{scope}[on background layer]
    \node[draw,dashed,rounded corners,fill=gray!3,inner sep=6pt,
          fit=(m1)(m3)(kms)(agent)(app),
          label={[font=\scriptsize\itshape]below:nodo SD-QKD (perimetro sicuro)}] (node) {};
  \end{scope}
  % nodi remoti
  \node[small] (r1) at (4.6,2.6) {nodo\\B};
  \node[small] (r2) at (4.4,-0.2) {nodo\\C};
  \node[small] (r3) at (-5.2,0.6) {nodo\\D};
  % collegamenti
  \draw[thick] (m3.east) -- node[lbl,above]{PHYS} (r2);
  \draw[thick] (m1.west) -- (r3);
  \draw[thick] (r1) -- (r2);
  \draw[thick,dotted,red!70!black] (kms.east) -- node[pos=0.55,font=\scriptsize,fill=white,inner sep=1pt]{VIRT via C} (r1);
  \draw[<->] (m1) -- (m1|-kms.south);
  \draw[<->] (m2) -- (m2|-kms.south);
  \draw[<->] (m3) -- (m3|-kms.south);
  \draw[<->] (agent.south) -- (agent.south|-kms.north);
  \draw[<->] (app.south) -- (app.south|-kms.north);
  \draw[<->,dashed,blue!60!black] (agent.north) -- node[pos=0.5,left,font=\scriptsize]{GS QKD 015} (ctl);
  \draw[<->,dashed,blue!60!black] (r1) -- (ctl);
  \draw[<->,dashed,blue!60!black] (r3) -- (ctl);
\end{tikzpicture}
\caption[Rete SD-QKD e struttura interna di un nodo]{Una rete SD-QKD: i nodi sono collegati da link fisici (linee
continue) e virtuali (punteggiata) e ciascuno è controllato dal controller
tramite il proprio SDN agent (tratteggio). Il nodo in primo piano mostra le
componenti interne: i moduli QKD, esposti come interfacce, il key manager,
l'interfaccia verso le applicazioni e l'agent. Adattata
da~\cite[Fig.~1]{qkd015}.}
\label{fig:sdqkd}
\end{figure}


\paragraph{4. Key association link.}
Tra due nodi la grandezza interessante non è il canale quantistico in sé
ma il fatto che i due nodi condividano chiave. La 015 modella proprio
questo.

\begin{definition}[Key association link]
Un \emph{key association link} (\texttt{qkd\_link}) è un'associazione
logica di chiave tra due nodi SD-QKD remoti. È di tipo \texttt{PHYS}
(diretto) se esiste un canale quantistico tra i due nodi, di tipo
\texttt{VIRT} (virtuale) se la chiave è ottenuta tramite key relay
attraverso nodi fidati intermedi.
\end{definition}

Un link fisico corrisponde a un QKD link della terminologia ITU-T; un link
virtuale corrisponde invece a una catena di KM link sopra più QKD link. Due
dettagli rendono il modello coerente. Il primo: un link virtuale
\emph{consuma} chiave dai link fisici su cui poggia, e per questo viene
registrato come \emph{applicazione interna} con una banda riservata
(\texttt{virt\_bandwidth}). Il secondo: ogni nodo intermedio di un link
virtuale conosce il nodo precedente e il successivo
(\texttt{virt\_prev\_hop}, \texttt{virt\_next\_hop}), quindi il cammino di
relay è configurato hop per hop dal controller. Perché un nodo possa fare
da relay deve dichiararlo con la capability
\texttt{key\_relay\_mode\_enable}.

\paragraph{5. Applicazioni.}
Dal punto di vista del nodo, un'applicazione è qualunque entità che chiede
chiavi al key manager.

\begin{definition}[Applicazione QKD]
Un'\emph{applicazione QKD} (\texttt{qkd\_app}) è un'istanza in esecuzione
che consuma chiavi da un nodo SD-QKD. È \emph{esterna}
(\texttt{CLIENT}) se si tratta di un utente, un HSM, un cifratore, una
funzione di rete virtuale o un protocollo di sicurezza; è \emph{interna}
(\texttt{INTERNAL}) se usa le chiavi per il funzionamento della rete
stessa, come l'autenticazione o il relay di un link virtuale.
\end{definition}

Ogni applicazione dichiara una qualità del servizio: banda massima e
minima di chiave (\texttt{max\_bandwidth}, \texttt{min\_bandwidth}, in
bit/s), \texttt{jitter} massimo della consegna in ms, tempo di vita
\texttt{ttl} di una chiave inutilizzata nel key store e la possibilità di
condividere cammini o chiavi tra più client. È la risposta tecnica
all'esigenza di QoS differenziata del caso UC2.

\paragraph{6. Controller SDN.}
Al vertice c'è il controller (nella 015 anche \emph{SD-QNC},
\emph{Software-Defined QKD Network Controller}), che ha la vista globale del
dominio e svolge tre ruoli. Fa da \emph{registro centrale} delle
applicazioni, così che un'applicazione debba conoscere solo il proprio key
manager e non la topologia della rete. Pianifica e crea i \emph{link},
fisici o virtuali, in base a domande e statistiche. Coordina infine il
proprio dominio con il controller della rete ottica (\emph{SD-ONC}) quando
serve un nuovo cammino ottico per un canale quantistico.

\paragraph{Che cosa non sale al controller.}
L'informazione che attraversa il piano di controllo è volutamente minima:
identificativi, stati, tassi, parametri fisici. Il materiale di chiave
\emph{non lascia mai il perimetro del nodo}. Per questo, osserva la
specifica, introdurre SDN in una rete a nodi fidati non aggiunge rischi
nuovi oltre a quelli già noti, come il denial of service.

% ---------------------------------------------------------------------
\subsubsection{Anatomia di un modello YANG}
\label{sec:yang-anatomia}

Sia la 015 sia la 018 sono scritte in YANG~\cite{rfc7950}, un linguaggio
nato in IETF per descrivere i dati di configurazione e di stato degli
apparati di rete e le operazioni che si possono invocare su di essi. Un
modello YANG non è codice eseguibile: è uno \emph{schema}, come un XSD per
XML o uno schema JSON, che stabilisce quali dati esistono, di che tipo sono,
come sono annidati e chi può modificarli. Su quello schema poi lavorano i
protocolli NETCONF e RESTCONF. Vediamo i costrutti principali, usando
esempi tratti dai file della 015~\cite{qkd015yang}.

\paragraph{Il modulo e la sua intestazione.}
L'unità di base è il \emph{modulo}. L'intestazione dichiara la versione
del linguaggio, un \texttt{namespace} globale univoco, un \texttt{prefix}
breve con cui gli altri moduli lo citano, i moduli importati e le revisioni.
\begin{lstlisting}[language=yang]
module etsi-qkd-sdn-node {
  yang-version "1";
  namespace "urn:etsi:qkd:yang:etsi-qkd-node";
  prefix "etsi-qkdn";
  import ietf-yang-types   { prefix "yang"; }
  import ietf-inet-types   { prefix "inet"; }
  import etsi-qkd-node-types { prefix "etsi-qkdn-types"; }
  organization "ETSI ISG QKD";
  revision "2022-01-30" { reference "ETSI GS QKD 015 V2.1.1"; }
  ...
}
\end{lstlisting}
La 015 è divisa in due moduli: \texttt{etsi-qkd-node-types} contiene solo
tipi, \texttt{etsi-qkd-sdn-node} contiene il modello vero e proprio. È una
scelta comune: separare i tipi permette ad altri moduli, come quelli della
018, di riusarli senza importare tutto il resto.

\paragraph{Tipi: built-in, typedef e identity.}
YANG offre tipi predefiniti (\texttt{boolean}, \texttt{string},
\texttt{uint32}, \texttt{decimal64}, \texttt{enumeration}\dots) e tipi
comuni definiti dall'IETF, come \texttt{yang:uuid} o
\texttt{yang:date-and-time}. Con \texttt{typedef} si definiscono tipi
derivati con un nome proprio. Per valori simbolici estendibili si usano le
\texttt{identity}: si dichiara un'identità base e poi identità che la
derivano; un campo di tipo \texttt{identityref} accetta qualunque identità
derivata dalla base. È così che la 015 definisce i tipi di link:
\begin{lstlisting}[language=yang]
identity QKD-LINK-TYPES { description "QKD key association link types."; }
identity VIRT { base QKD-LINK-TYPES; description "Virtual Link."; }
identity PHYS { base QKD-LINK-TYPES; description "Physical Link."; }

typedef qkd-link-types {
  type identityref { base QKD-LINK-TYPES; }
}
\end{lstlisting}
Il vantaggio rispetto a un'\texttt{enumeration} è che un altro modulo può
aggiungere, per esempio, un nuovo tipo di link senza modificare l'originale.
Allo stesso modo sono definiti le tecnologie (\texttt{CV-QKD},
\texttt{DV-QKD}\dots), gli stati di nodo (\texttt{NEW},
\texttt{OPERATING}, \texttt{DOWN}, \texttt{FAILURE}, \texttt{OUT}), i ruoli
(\texttt{TRANSMITTER}, \texttt{RECEIVER}, \texttt{TRANSCEIVER}) e i tipi
di applicazione (\texttt{CLIENT}, \texttt{INTERNAL}).

\paragraph{Nodi dati: container, list, leaf, leaf-list.}
I dati sono organizzati ad albero con quattro costrutti:
\begin{itemize}
  \item \texttt{leaf}: un singolo valore tipizzato (per esempio
        \texttt{qkdn\_id});
  \item \texttt{leaf-list}: una sequenza di valori dello stesso tipo
        (per esempio \texttt{virt\_next\_hop}, che ammette più successori);
  \item \texttt{container}: un nodo che raggruppa altri nodi, senza
        valore proprio (per esempio \texttt{qkdn\_capabilities});
  \item \texttt{list}: una collezione di elementi strutturati, ciascuno
        identificato da una \texttt{key} (per esempio la lista
        \texttt{qkd\_link} con chiave \texttt{qkdl\_id}).
\end{itemize}

\paragraph{Configurazione e stato.}
Ogni nodo è di \emph{configurazione} (\texttt{config true}, il default:
il client può scriverlo) oppure di \emph{stato operativo}
(\texttt{config false}: lo produce l'apparato e il client può solo
leggerlo). Nella 015, per esempio, \texttt{qkdl\_enable} è configurabile,
mentre \texttt{qkdn\_status}, \texttt{skr} ed \texttt{eskr} sono stato.
Questa distinzione riflette la divisione tra ciò che il controller
\emph{decide} e ciò che il nodo \emph{misura}.

\paragraph{Riuso ed estensione: grouping, uses, augment.}
Un \texttt{grouping} è un blocco di nodi riutilizzabile, che viene
copiato dove compare \texttt{uses}. Il modello della 015 è costruito quasi
tutto così: \texttt{app\_basic}, \texttt{app\_details}, \texttt{app\_qos},
\texttt{qkd\_links} e così via sono grouping poi assemblati nel container
\texttt{qkd\_node}. L'\texttt{augment} fa l'operazione opposta: aggiunge
nodi a un albero definito \emph{in un altro modulo}, senza toccarlo. È lo
strumento che la 018 userà per arricchire la propria topologia con i dati
di inventario.

\paragraph{Riferimenti e vincoli.}
Un \texttt{leafref} è un campo il cui valore deve coincidere con un valore
esistente altrove nell'albero, come una chiave esterna in un database. Le
istruzioni \texttt{when} e \texttt{must} aggiungono condizioni, per esempio
che un campo esista solo se il link è di un certo tipo. La 015 indica
proprio così che i campi \texttt{virt\_*} valgono solo per link virtuali e
i campi \texttt{phys\_*} solo per quelli fisici.

\paragraph{Operazioni e notifiche: rpc e notification.}
Oltre ai dati, un modulo può definire operazioni invocabili dal client
(\texttt{rpc}, con \texttt{input} e \texttt{output}) e messaggi che il
server invia spontaneamente (\texttt{notification}). La 015 definisce
soprattutto notifiche; la 018, come vedremo, definisce anche RPC.

\paragraph{Leggere un albero.}
Lo strumento \texttt{pyang} produce una vista ad albero di un modulo, che
useremo nel seguito. La notazione è compatta: \texttt{rw} indica un nodo di
configurazione, \texttt{ro} un nodo di stato, \texttt{?} un nodo
facoltativo, \texttt{*} una lista o leaf-list, la chiave di una lista
compare tra parentesi quadre e \texttt{-{}-{}-n} indica una notifica.

\paragraph{Dai dati al messaggio.}
Lo stesso schema produce sia XML sia JSON. In JSON~\cite{rfc7951} il primo
livello è qualificato con il nome del modulo e i valori di tipo
\texttt{identityref} riportano il modulo che definisce l'identità. Un link
fisico esposto da un nodo apparirebbe così:
\begin{lstlisting}
{ "etsi-qkd-sdn-node:qkd_node": {
    "qkdn_id": "0f3c...-a1",
    "qkd_links": { "qkd_link": [ {
        "qkdl_id": "7b2e...-c4",
        "qkdl_type": "etsi-qkd-node-types:PHYS",
        "qkdl_local":  { "qkdn_id": "0f3c...-a1", "qkdi_id": 1 },
        "qkdl_remote": { "qkdn_id": "9d41...-e7", "qkdi_id": 2 },
        "phys_wavelength": 1550,
        "qkdl_performance": { "skr": 20000, "eskr": 12000,
                              "expected_consumption": 8000 } } ] } } }
\end{lstlisting}
(I valori sono illustrativi.)

% ---------------------------------------------------------------------
\subsubsection{Il modello del nodo SD-QKD}

Con questi strumenti il modello della 015 si legge facilmente. Il container
radice \texttt{qkd\_node} ha quattro sezioni, una per ciascuna componente
descritta sopra, più le notifiche:
\begin{lstlisting}
module: etsi-qkd-sdn-node
  +--rw qkd_node
     +--rw qkdn_id              yang:uuid
     +--ro qkdn_status?         etsi-qkdn-types:qkdn-status-types
     +--rw qkdn_version?        string
     +--rw qkdn_location_id?    string
     +--rw qkdn_capabilities
     |  +--rw link_stats_support?          boolean
     |  +--rw application_stats_support?   boolean
     |  +--rw key_relay_mode_enable?       boolean
     +--rw qkd_applications
     |  +--rw qkd_app* [app_id]
     |     +--rw app_id, app_type, server_app_id, client_app_id*, ...
     |     +--rw app_qos            (max/min_bandwidth, jitter, ttl, ...)
     |     +--rw backing_qkdl_id*, local_qkdn_id, remote_qkdn_id
     +--rw qkd_interfaces
     |  +--rw qkd_interface* [qkdi_id]
     |     +--rw qkdi_id, qkdi_model, qkdi_type
     |     +--rw qkdi_att_point     (device, port)
     |     +--rw qkdi_capabilities  (role, wavelength_range, max_absorption)
     |     +--ro qkdi_status
     +--rw qkd_links
        +--rw qkd_link* [qkdl_id]
           +--rw qkdl_id, qkdl_status, qkdl_enable, qkdl_type
           +--rw qkdl_local / qkdl_remote   (qkdn_id, qkdi_id)
           +--rw qkdl_applications*
           +--rw virt_prev_hop, virt_next_hop*, virt_bandwidth   (VIRT)
           +--rw phys_channel_att, phys_wavelength, phys_qkd_role (PHYS)
           +--rw qkdl_performance   (expected_consumption, skr, eskr,
                                     phys_perf* [perf_type]: QBER, SNR)
\end{lstlisting}

\paragraph{Identità del nodo.}
\texttt{qkdn\_id} è un UUID univoco nella rete; \texttt{qkdn\_status},
\texttt{qkdn\_version} e \texttt{qkdn\_location\_id} servono a inventario e
diagnostica.

\paragraph{Capability.}
Le tre capability dicono al controller che cosa il nodo è disposto a fare:
esporre statistiche sui link (SKR, consumo, stato, QBER), esporre
statistiche sulle applicazioni e fare da relay. Un nodo con
\texttt{key\_relay\_mode\_enable = false} non verrà mai scelto come nodo
intermedio di un link virtuale.

\paragraph{Applicazioni.}
Ogni \texttt{qkd\_app} lega un'applicazione al servizio di chiave:
\texttt{server\_app\_id} identifica l'entità che ha avviato la richiesta,
\texttt{client\_app\_id} le entità autorizzate a ricevere la stessa chiave,
\texttt{backing\_qkdl\_id} il link che fornisce le chiavi,
\texttt{local\_qkdn\_id} e \texttt{remote\_qkdn\_id} i due nodi coinvolti.
Finché il nodo remoto non è noto, il nodo locale non può servire chiavi.
Le statistiche (\texttt{app\_statistics}) riportano i bit consumati per
intervallo di tempo.

\paragraph{Interfacce.}
Ogni \texttt{qkd\_interface} descrive un modulo: modello, tecnologia, stato,
punto di attacco a un dispositivo ottico (\texttt{qkdi\_att\_point}) e
capability fisiche. È l'informazione che serve al controller per abbinare
due interfacce compatibili e chiedere alla rete ottica un cammino tra le
rispettive porte.

\paragraph{Link e prestazioni.}
Ogni \texttt{qkd\_link} riporta gli estremi, il tipo, lo stato, le
applicazioni servite e i parametri specifici del tipo. Il blocco
\texttt{qkdl\_performance} contiene le tre grandezze che guidano tutta la
pianificazione:
\begin{itemize}
  \item \texttt{skr}: tasso di generazione di chiave segreta, in bit/s;
  \item \texttt{expected\_consumption}: somma delle bande di tutte le
        applicazioni sul link;
  \item \texttt{eskr}: tasso \emph{effettivo}, cioè disponibile dopo il
        consumo interno.
\end{itemize}
In forma compatta, se $\mathcal{A}_\ell$ è l'insieme delle applicazioni sul
link $\ell$, ciascuna con banda $b_a$,
\begin{equation}
  \text{expected\_consumption}_\ell = \sum_{a\in\mathcal{A}_\ell} b_a ,
  \qquad
  \mathrm{ESKR}_\ell \approx \mathrm{SKR}_\ell - \text{expected\_consumption}_\ell .
  \label{eq:eskr015}
\end{equation}
La seconda relazione è una lettura intuitiva dei campi, non una formula
dello standard. Per i link fisici si aggiungono metriche come QBER e SNR.

% ---------------------------------------------------------------------
\subsubsection{Notifiche}

Interrogare continuamente i nodi sarebbe inefficiente. La 015 prevede
quindi notifiche YANG dal nodo al controller, divise per oggetto:
\begin{itemize}
  \item \textbf{applicazioni}: \texttt{sdqkdn\_application\_new} (una nuova
        applicazione chiede chiavi), \texttt{sdqkdn\_application\_qos\_update},
        \texttt{sdqkdn\_application\_disconnected};
  \item \textbf{interfacce}: \texttt{sdqkdn\_interface\_new},
        \texttt{sdqkdn\_interface\_down} (con un campo \emph{reason}),
        \texttt{sdqkdn\_interface\_out};
  \item \textbf{link}: \texttt{sdqkdn\_link\_down},
        \texttt{sdqkdn\_link\_perf\_update} (variazione significativa delle
        prestazioni), \texttt{sdqkdn\_link\_overloaded} (il link non regge
        la domanda).
\end{itemize}
A queste si aggiunge una struttura generica di eventi e allarmi, pensata
per accogliere nuovi tipi di evento senza cambiare il modello. Inoltre i
protocolli basati su YANG permettono di sottoscrivere qualunque dato del
datastore, così che ogni variazione possa essere segnalata in tempo reale.

% ---------------------------------------------------------------------
\subsubsection{I workflow}

La clausola 5 descrive tre scenari fondamentali, che mostrano le componenti
all'opera.

\paragraph{Registrazione di un'applicazione.}
È il caso più istruttivo, perché spiega perché il controller fa da registro
centrale (Fig.~\ref{fig:seq015}). Un'applicazione conosce solo l'indirizzo
del proprio key manager; non conosce il nodo dell'applicazione con cui
vuole comunicare.
\begin{enumerate}
  \item L'applicazione nel dominio di sicurezza $A$ si collega al proprio
        nodo tramite la key delivery API, inviando il proprio ID, l'ID
        dell'applicazione remota e la QoS richiesta.
  \item Il nodo $A$ non sa dove si trovi l'applicazione remota e inoltra la
        richiesta al controller, aggiungendo il proprio ID.
  \item Se l'applicazione remota non è ancora registrata, il controller
        risponde con un semplice riscontro.
  \item Nel dominio $B$ l'applicazione remota si collega al proprio nodo, che
        a sua volta informa il controller.
  \item Il controller riconosce la coppia, genera un identificativo unico
        di servizio (\texttt{app\_id}) e invia a entrambi i nodi tutte le
        informazioni necessarie: ID di applicazioni e nodi, \texttt{app\_id}
        e QoS.
  \item I nodi comunicano l'\texttt{app\_id} alle applicazioni e riservano
        chiavi per esse; da qui in poi le applicazioni chiedono chiavi
        usando quell'identificativo.
\end{enumerate}

% Figura: registrazione di un'applicazione (adattata da GS QKD 015, Fig. 2)
\begin{figure}[t]
\centering
\begin{tikzpicture}[font=\scriptsize]
  \def\xA{0}\def\xB{3.0}\def\xC{6.6}\def\xD{10.2}\def\xE{13.2}
  \node[actor,minimum width=2cm] at (\xA,0) {App A};
  \node[actor,minimum width=2cm,fill=green!8] at (\xB,0) {Nodo A};
  \node[actor,minimum width=2cm,fill=blue!8] at (\xC,0) {Controller};
  \node[actor,minimum width=2cm,fill=green!8] at (\xD,0) {Nodo B};
  \node[actor,minimum width=2cm] at (\xE,0) {App B};
  \foreach \x in {\xA,\xB,\xC,\xD,\xE} \draw[gray] (\x,-0.4) -- (\x,-6.6);
  \draw[msg] (\xA,-0.9) -- node[above]{ID, ID peer, QoS} (\xB,-0.9);
  \draw[msg] (\xB,-1.5) -- node[above]{\texttt{sdqkdn\_application\_new}} (\xC,-1.5);
  \draw[ret] (\xC,-2.1) -- node[above]{ack (peer ignoto)} (\xB,-2.1);
  \draw[msg] (\xE,-2.9) -- node[above]{ID, ID peer, QoS} (\xD,-2.9);
  \draw[msg] (\xD,-3.5) -- node[above]{\texttt{sdqkdn\_application\_new}} (\xC,-3.5);
  \draw[msg] (\xC,-4.1) -- ++(0.35,0) -- ++(0,-0.35) -- (\xC,-4.45);
  \node[right] at (\xC+0.4,-4.27) {abbina e crea \texttt{app\_id}};
  \draw[msg] (\xC,-5.1) -- node[above]{config.\ (\texttt{app\_id}, QoS)} (\xB,-5.1);
  \draw[msg] (\xC,-5.1) -- node[above]{config.\ (\texttt{app\_id}, QoS)} (\xD,-5.1);
  \draw[ret] (\xB,-5.9) -- node[above]{\texttt{app\_id}} (\xA,-5.9);
  \draw[ret] (\xD,-5.9) -- node[above]{\texttt{app\_id}} (\xE,-5.9);
\end{tikzpicture}
\caption[Registrazione di un'applicazione in una rete SD-QKD]{Registrazione di un'applicazione in una rete SD-QKD. Le
applicazioni conoscono solo il proprio key manager; è il controller ad
abbinare le due richieste e a creare l'identificativo di servizio.
Adattata da~\cite[Fig.~2]{qkd015}.}
\label{fig:seq015}
\end{figure}


\paragraph{Creazione di un link fisico.}
Il controller riceve un innesco, da una domanda rilevata sui nodi o dalla
propria NBI. Cerca nei due siti due interfacce libere e compatibili,
chiede al controller della rete ottica (SD-ONC) un cammino ottico
(\emph{lightpath}) con i vincoli del canale quantistico e, ricevuta la
conferma, configura le due interfacce sui rispettivi nodi. Se i due
controller sono integrati in uno solo, l'operazione è interna. Questo
workflow è il cuore della coesistenza tra QKD e rete ottica chiesta dal
caso UC4.

\paragraph{Creazione di un link virtuale.}
Se non ci sono interfacce compatibili o non esiste un lightpath adatto, il
controller ripiega su un link virtuale. Calcola un cammino sulla rete QKD
esistente, per esempio su un grafo pesato costruito con disponibilità di
chiave, tassi di generazione e consumi; configura ogni nodo intermedio in
modalità relay indicando precedente e successivo; infine configura gli
estremi con la banda da fornire. La specifica non prescrive l'algoritmo.
Una scelta naturale, coerente con la~\eqref{eq:eskr015}, è pesare ciascun
link con l'inverso della chiave disponibile e scartare i link che non
reggono la banda $r$ richiesta:
\begin{equation}
  w_\ell = \frac{1}{\mathrm{ESKR}_\ell}, \qquad
  p^\star = \arg\min_{p\,:\,\min_{\ell\in p}\mathrm{ESKR}_\ell \ge r}
            \;\sum_{\ell\in p} w_\ell .
  \label{eq:path015}
\end{equation}
La~\eqref{eq:path015} è un esempio nostro, non parte dello standard. La
specifica avverte che, se gestiti male, i link virtuali possono congestionare
i link fisici sottostanti: è proprio per questo che serve un controller con
la vista globale di consumi e disponibilità.

% ---------------------------------------------------------------------
\subsubsection{Protocolli e sicurezza}

La 015 non impone un protocollo, ma indica i due più diffusi tra quelli
basati su YANG: NETCONF~\cite{rfc6241}, tipicamente tra apparato e
controller, transazionale e trasportato su SSH o TLS con codifica XML; e
RESTCONF~\cite{rfc8040}, tipicamente sulla NBI dei controller, basato su
HTTP con operazioni CRUD e codifica JSON o XML.

Sul piano della sicurezza (clausola 6) la specifica distingue due casi
all'installazione di un nuovo sistema. Se il sistema entra in un'area
sicura già presente nella rete, il canale verso il controller può essere
protetto con chiavi QKD prodotte da link esistenti. Se invece si aggiunge
un nuovo luogo sicuro, serve uno schema iniziale (un'identità fidata o un
segreto pre-condiviso, usato solo durante l'installazione) per presentare
il nuovo SDN agent al controller; dopo la prima autenticazione si passa a
chiavi QKD. È lo stesso principio di key growing visto in apertura.

% ---------------------------------------------------------------------
\subsubsection{Dal dominio singolo all'orchestrazione}

La 015 risolve il controllo di \emph{un} dominio QKD da parte di \emph{un}
controller. Restano fuori tre problemi che la GS~QKD~002 aveva già fatto
intravedere. Primo, un operatore può avere più domini QKD, magari di vendor
diversi, ciascuno con il proprio controller. Secondo, le applicazioni che
consumano chiave possono stare in un'altra rete, per esempio nei nodi della
rete di trasporto ottico, controllata da un controller diverso. Terzo,
l'operatore vuole gestire la QKD con gli stessi strumenti di orchestrazione
che usa per il resto della rete. Serve quindi un livello superiore e
un'interfaccia standard per raggiungerlo: è il compito della GS~QKD~018.

% =====================================================================
\section{ETSI GS QKD 018}
\label{sec:qkd018}
% =====================================================================

La GS~QKD~018 (V1.1.1, aprile 2022)~\cite{qkd018} definisce
l'\emph{orchestration interface} tra un SDN orchestrator e l'SDN controller
di una rete QKD. Il controller agisce da \emph{server}, l'orchestrator da
\emph{client}; il modello informativo è espresso in YANG ed è indipendente
dal vendor. Nella terminologia SDN è la \emph{northbound interface} del
controller QKD, cioè l'esatto complemento della GS~QKD~015. La gerarchia
complessiva diventa:
\begin{center}
\small
orchestrator $\xrightarrow{\text{QKD 018}}$ controller
$\xrightarrow{\text{QKD 015}}$ SDN agent / nodo SD-QKD
$\xrightarrow{\text{QKD 004/014}}$ SAE.
\end{center}

\begin{definition}[Orchestrazione SDN]
L'\emph{orchestrazione SDN} è il processo continuo con cui le risorse
disponibili vengono coordinate automaticamente, secondo criteri di
ottimizzazione, per stabilire e rilasciare servizi end-to-end che
attraversano più domini di rete, ciascuno governato dal proprio controller
SDN~\cite[§4.2]{qkd018}.
\end{definition}

\begin{definition}[Service link]
Un \emph{service link} è un collegamento logico di associazione di chiave
tra due nodi QKD, realizzato da uno o più link fisici QKD (single-hop o
multi-hop). A livello di orchestrazione corrisponde al link virtuale della
GS~QKD~015.
\end{definition}

% Figura: scenario di riferimento della GS QKD 018 (adattato da GS QKD 018, Fig. 1)
\begin{figure}[t]
\centering
\begin{tikzpicture}
  \node[box,fill=blue!8,minimum width=3.3cm] (orch) {SDN orchestrator};
  \node[box,fill=green!8,minimum width=3.3cm,below left=12mm and 6mm of orch] (cq)
       {SDN controller\\rete QKD};
  \node[box,fill=orange!10,minimum width=3.3cm,below right=12mm and 6mm of orch] (co)
       {SDN controller\\rete OTN};
  \node[box,fill=green!4,below=10mm of cq,minimum width=4.6cm] (nq)
       {nodi SD-QKD (agent, KME)};
  \node[box,fill=orange!5,below=10mm of co,minimum width=4.6cm] (no)
       {nodi OTN (SAE)};
  \draw[<->,very thick,red!70!black] (orch) -- node[lbl]{GS QKD 018} (cq);
  \draw[<->,thick] (orch) -- node[lbl]{NBI OTN (fuori scope)} (co);
  \draw[<->,thick] (cq) -- node[lbl]{GS QKD 015} (nq);
  \draw[<->,thick] (co) -- (no);
  \draw[<->,dashed] (nq) -- node[lbl]{QKD 004 / 014} (no);
\end{tikzpicture}
\caption[Scenario di riferimento della GS QKD 018]{Scenario di riferimento della GS~QKD~018 (adattato
da~\cite[Fig.~1]{qkd018}). In rosso l'interfaccia standardizzata; la
linea tratteggiata è la consegna delle chiavi dai KME ai SAE, che avviene
nello stesso perimetro di sicurezza.}
\label{fig:scenario}
\end{figure}


% ---------------------------------------------------------------------
\subsection{Use case della QKD 018}
\label{sec:usecase018}

Il caso d'uso da cui nasce lo standard (§4.1) riprende quasi alla lettera
il problema del backbone della GS~QKD~002. Un operatore che introduce la
QKD può far convivere canali quantistici e classici nella stessa fibra
oppure tenerli separati. La convivenza è difficile per due motivi: le
potenze ottiche dei canali classici sono di molti ordini di grandezza
superiori e disturbano il canale quantistico, e le distanze raggiungibili
dai due tipi di canale sono molto diverse. In un ambiente ricco di fibra
l'operatore può quindi scegliere una rete QKD separata dalla rete di
trasporto ottico (OTN), con vantaggi di prestazioni, isolamento e
separazione delle responsabilità. Seguendo il paradigma SDN, le due reti
sono governate da controller diversi.

Nasce così un problema di coordinamento: i SAE che devono cifrare il
traffico stanno nei nodi OTN, le chiavi sono generate nei nodi QKD e
ciascun controller conosce solo la propria rete. L'orchestrator risolve tre
questioni:

\begin{enumerate}
  \item \textbf{Address matching}: associa l'indirizzo del nodo OTN che deve
        ricevere chiavi a quello del nodo QKD che le fornisce; i due si
        trovano nello stesso perimetro di sicurezza (Fig.~\ref{fig:scenario});
  \item \textbf{provisioning end-to-end}: traduce la richiesta di un cliente
        nella configurazione dei singoli domini, sceglie i domini da
        attraversare e coordina i cammini;
  \item \textbf{gestione}: raccoglie topologia, inventario delle risorse di
        chiave e informazioni FCAPS di ogni dominio.
\end{enumerate}

Generalizzando, lo stesso schema permette di orchestrare \emph{più domini
QKD di vendor diversi}, ciascuno dietro il proprio controller, oltre a
domini QKD e classici insieme. Lo standard sottolinea anche un vantaggio
pratico: un operatore che ha già un orchestrator per le proprie reti
ottiche può integrare la rete QKD senza introdurre strumenti di gestione
separati, che è esattamente il requisito di gestione integrata del caso UC4.

Le funzioni richieste all'orchestrator (§4.2) sono:
\begin{itemize}
  \item tradurre le richieste di servizio dei clienti nella configurazione
        dei domini e nell'allocazione dei SAE;
  \item richiedere attivazione e rilascio dei servizi, incluse le
        connessioni tra domini;
  \item calcolare cammini multi-dominio, con le connessioni inter-dominio e
        gli estremi;
  \item modificare i cammini su vincoli dell'operatore;
  \item scoprire nodi e link di ciascun controller e mantenerne una vista
        astratta;
  \item monitorare le risorse di chiave disponibili su nodi e link;
  \item monitorare la gestione FCAPS e riceverne le notifiche.
\end{itemize}

% ---------------------------------------------------------------------
\subsection{Il modello informativo}
\label{sec:infomodel}

La clausola 5 traduce queste funzioni in cinque blocchi informativi.

\paragraph{Discovery della topologia (§5.1).}
Il controller espone nodi SD-QKD, interfacce e link fisici, e in seguito i
service link. L'orchestrator non ha bisogno della composizione fisica
completa, ma di una vista \emph{astratta} di ciascun dominio e della
connettività tra domini. La discovery può essere \emph{proattiva}
(l'orchestrator interroga il controller prima di ogni nuovo calcolo di
cammino) o \emph{reattiva} (il controller notifica ogni cambiamento, per
esempio l'aggiunta o la rimozione di un nodo).

\paragraph{Stato della rete e inventario (§5.2).}
Per ogni nodo e link fisico il controller espone lo stato, le prestazioni
fisiche e soprattutto le risorse di chiave: il tasso di generazione e il
tasso disponibile dopo il consumo interno. Per i service link si aggiungono
la banda, la chiave disponibile e la politica di ripristino in caso di
guasto.

\paragraph{Stato dei servizi end-to-end (§5.3).}
L'orchestrator può interrogare le applicazioni esterne attive, i nodi
locale e remoto e i service link che le servono, con il relativo stato e
tasso di chiave.

\paragraph{Provisioning con calcolo dei cammini (§5.4).}
L'orchestrator chiede al controller cammini che soddisfino certi vincoli; il
controller restituisce una lista di candidati, eventualmente più volte con
vincoli diversi, finché l'orchestrator non ne sceglie uno da attivare. In
caso di guasto l'orchestrator può chiedere di sostituire il cammino.

\paragraph{Notifiche (§5.5).}
Il controller segnala eventi, allarmi e superamenti di soglia. Le soglie
sulle prestazioni (QBER, SKR, ESKR) sono configurabili dall'orchestrator.
Rispetto alla 015 si aggiungono notifiche sullo stato dei nodi
(\texttt{sdqkdn\_new}, \texttt{\_operating}, \texttt{\_down},
\texttt{\_failure}, \texttt{\_out}), sulle applicazioni attivate
(\texttt{sdqkdn\_application\_deployed}) e sui service link
(\texttt{sdqkdn\_svc\_link\_down}, \texttt{\_perf\_update},
\texttt{\_overloaded}).

La~\eqref{eq:eskr015} si estende in modo naturale al livello di
orchestrazione. Poiché in un cammino multi-hop $p$ ogni link inoltra la
stessa chiave, un nuovo servizio con richiesta $r$ è sostenibile solo se il
collo di bottiglia lo regge:
\begin{equation}
  p \text{ ammissibile} \iff
  \min_{\ell\in p}\, \mathrm{ESKR}_\ell \;\ge\; r .
  \label{eq:admissible}
\end{equation}
Anche questa è una formalizzazione nostra, coerente con i campi
\texttt{skr}, \texttt{eskr} ed \texttt{expected\_consumption} dei modelli
YANG; lo standard non prescrive alcun algoritmo di calcolo dei cammini.

% =====================================================================
\subsection{I moduli YANG}
\label{sec:yang018}
% =====================================================================

L'Annex~A (normativo) traduce il modello in tre moduli YANG~1.1, pubblicati
su ETSI Forge~\cite{qkd018yang}. Con i costrutti visti in
§\ref{sec:yang-anatomia} la loro struttura si legge rapidamente. I moduli
importano \texttt{etsi-qkd-sdn-node} ed \texttt{etsi-qkd-node-types} della
GS~QKD~015 per riusarne i tipi (identificativi, tipi di link, QoS delle
applicazioni): la 018 non reinventa il nodo, ne fornisce una vista
astratta. I tre moduli formano una catena: \emph{topology} (lo scheletro)
$\to$ \emph{inventory} (che cosa c'è e in che stato) $\to$
\emph{connectivity} (che cosa ci si fa).

\paragraph{\texttt{etsi-qkd-sdn-topology} (prefisso \texttt{qkdnw-tplg}).}
Descrive solo la struttura, cioè chi è collegato a chi:
\begin{lstlisting}
qkd_network
 +-- sdqkd_nodes
 |    +-- qkdn* [qkdn_id]                          nodi SD-QKD
 |         +-- qkd_interfaces/qkdi* [qkdi_id]      moduli QKD
 +-- qkd_links
      +-- qkd_phys_links/phys_link* [phys_link_id]  link fisici
      |    +-- link_type
      |    +-- phys_link_local_qkdn   (qkdn_id_ref, qkdi_id_ref)
      |    +-- phys_link_remote_qkdn  (qkdn_id_ref, qkdi_id_ref)
      +-- qkd_svc_links/svc_link* [svc_link_id]     service link
           +-- ... estremi locale/remoto come sopra
\end{lstlisting}
Rispetto alla 015 il punto di vista cambia: non c'è più un albero per
nodo, ma un unico albero per l'intera rete del dominio. In termini astratti
è un grafo $G=(V,E)$ con $V$ insieme dei nodi \texttt{qkdn} ed
$E = E_{\text{phys}}\cup E_{\text{svc}}$, dove i \texttt{phys\_link}
rappresentano adiacenze fisiche e i \texttt{svc\_link} relazioni logiche di
chiave, anche multi-hop. Gli estremi sono \texttt{leafref}, cioè
riferimenti validati a nodi e interfacce esistenti. Avendo il grafo $G_i$ di
ciascun dominio $i$, l'orchestrator compone una vista multi-dominio
\begin{equation}
  G = \Big(\bigcup_i V_i,\; \bigcup_i E_i \,\cup\, E_{\text{inter}}\Big),
\end{equation}
dove $E_{\text{inter}}$ sono le adiacenze tra domini, tipicamente
realizzate da nodi di bordo condivisi.

\paragraph{\texttt{etsi-qkd-sdn-inventory} (prefisso \texttt{qkdnw-invt}).}
Non crea un nuovo albero ma fa \texttt{augment} di quello di topologia,
aggiungendo foglie agli elementi già definiti (Tab.~\ref{tab:inventory}).
È un esempio da manuale dell'uso di \texttt{augment}: chi è interessato
solo alla struttura implementa la topologia, chi vuole anche lo stato
aggiunge l'inventario senza duplicare nulla.

\begin{table}[h]
\centering\small
\caption{Elementi aggiunti dal modulo inventory all'albero di topologia.}
\label{tab:inventory}
\begin{tabularx}{\textwidth}{@{}l L@{}}
\toprule
\textbf{Elemento} & \textbf{Foglie aggiunte} \\
\midrule
\texttt{qkdn} & \texttt{qkdn\_version}, \texttt{qkdn\_capabilities},
\texttt{qkdn\_location\_id}, \texttt{KME\_ip\_port} (IP e porta del KME, dove
i SAE prelevano le chiavi via QKD~014), \texttt{qkdn\_status} \\
\texttt{qkdi} & \texttt{role\_support}, \texttt{qkdi\_model},
\texttt{qkdi\_type}, \texttt{qkdi\_att\_point}, \texttt{qkdi\_status} \\
\texttt{phys\_link} & \texttt{enable} (configurabile), \texttt{wavelength},
\texttt{performance} (\texttt{skr}, \texttt{eskr},
\texttt{expected\_consumption} e la lista \texttt{phys\_link\_perf} con QBER
e SNR), \texttt{phys\_link\_status} \\
\texttt{svc\_link} & \texttt{bandwidth}, \texttt{performance},
\texttt{path\_restoration} (\texttt{type} \texttt{AUTOMATIC}/\texttt{MANUAL},
\texttt{hold\_off\_time} in ms), \texttt{svc\_link\_status} \\
\bottomrule
\end{tabularx}
\end{table}

Il modulo definisce anche le notifiche che realizzano la parte reattiva
della §5.5: \texttt{sdqkdn\_new}, \texttt{sdqkdn\_interface\_new},
\texttt{sdqkdn\_phys\_link\_new}, \texttt{sdqkdn\_svc\_link\_new}, le
varianti \texttt{\_perf\_update} e \texttt{\_overloaded}, distinte per link
fisici e service link (ad esempio \texttt{sdqkdn\_phys\_link\_overloaded}),
e le generiche \texttt{alarm} ed \texttt{event}, che riportano stato,
messaggio e severità di nodi, interfacce e link. Le notifiche sullo stato del
nodo elencate nella §5.5 (\texttt{sdqkdn\_operating}, \texttt{\_down},
\texttt{\_failure}, \texttt{\_out}) non hanno una \texttt{notification}
dedicata nel modulo: lo stato viaggia nel campo \texttt{qkdn\_status} di
\texttt{alarm} ed \texttt{event}. Il campo \texttt{eskr} dei link è il dato
da cui l'orchestrator capisce se un cammino regge un nuovo servizio,
secondo la~\eqref{eq:admissible}.

\paragraph{\texttt{etsi-qkd-sdn-connectivity} (prefisso \texttt{qkdnw-conn}).}
È il modulo operativo: gestisce le applicazioni e i cammini che le servono.
\begin{lstlisting}
connectivity_service/qkd_service/qkd_app* [app_id]
 +-- (dati dell'applicazione da QKD 015: stato, tipo, QoS, statistiche)
 +-- app_mapped_svc_link
      +-- svc_link_and_qkdn_ids   svc_link_id_ref + estremi (nodo, interfaccia)
      +-- path_constraints
      |    +-- include_nodes* [qkdn_id_ref qkdi_id_ref]
      |    +-- exclude_nodes* [qkdn_id_ref qkdi_id_ref]
      +-- candidate_paths
      |    +-- num_of_candidate_paths
      |    +-- candidate_path* [path_id]   phys_links*  (path_id = rango)
      +-- deploy_path             path_id + phys_links* del cammino scelto
\end{lstlisting}
Qui compaiono le \texttt{rpc}, che ricalcano passo per passo il workflow di
provisioning:
\begin{description}[font=\normalfont\ttfamily\footnotesize,style=nextline,leftmargin=1.2em]
  \item[app-registration] registra l'applicazione nel dominio;
  \item[create-service-link-id-for-app] crea il service link tra nodo
        locale e remoto;
  \item[create-candidate-paths-per-service-link-for-app] il controller
        calcola i cammini candidati rispettando QoS e vincoli
        \texttt{include\_nodes}/\texttt{exclude\_nodes}, espressi come
        coppie (nodo, interfaccia);
  \item[update-candidate-paths-per-service-link-for-app] ricalcolo con
        vincoli diversi;
  \item[deploy-a-candidate-path-per-service-link-for-app] l'orchestrator
        sceglie un \texttt{path\_id} e il controller lo attiva;
  \item[delete-deployed-path-per-service-link-for-app] rilascio del
        cammino.
\end{description}
Le notifiche del modulo sono \texttt{sdqkdn\_application\_new},
\texttt{sdqkdn\_application\_deployed},
\texttt{sdqkdn\_application\_qos\_update} e
\texttt{sdqkdn\_application\_disconnected}.

Il punto concettuale è la \textbf{divisione dei ruoli}: il controller
conosce il proprio dominio e \emph{propone} i cammini, l'orchestrator ha la
vista d'insieme e \emph{decide}. La Tab.~\ref{tab:moduli} riassume.

\begin{table}[h]
\centering\small
\caption{I tre moduli YANG della GS~QKD~018.}
\label{tab:moduli}
\begin{tabularx}{\textwidth}{@{}l L l L@{}}
\toprule
\textbf{Modulo} & \textbf{Domanda} & \textbf{Clausola} & \textbf{Dati} \\
\midrule
topology & Com'è fatta la rete? & §5.1 & struttura (grafo) \\
inventory & Quante risorse ci sono e in che stato? & §5.2, §5.5 &
stato, prestazioni, configurazione, notifiche \\
connectivity & Come servo un'applicazione end-to-end? & §5.3, §5.4, §6.2 &
applicazioni, RPC di calcolo e deploy dei cammini \\
\bottomrule
\end{tabularx}
\end{table}

% =====================================================================
\subsection{Workflow di provisioning multi-dominio}
\label{sec:workflow}
% =====================================================================

La clausola 6 lega le informazioni al ciclo di vita del servizio. Quando la
rete QKD viene introdotta, l'orchestrator ne acquisisce la topologia
(§5.1); durante l'esercizio riceve, su richiesta o periodicamente, stato e
inventario (§5.2); dopo l'attivazione di un servizio ne monitora lo stato
(§5.3); in ogni momento riceve notifiche (§5.5).

La §6.2 descrive il caso completo (Fig.~\ref{fig:seq018}): un orchestrator
coordina il controller QKD e il controller OTN, che non comunicano tra loro
direttamente.

% Figura: provisioning end-to-end QKD + OTN (adattata da GS QKD 018, Fig. 2)
\begin{figure}[t]
\centering
\begin{tikzpicture}[font=\scriptsize]
  \def\xA{0}\def\xB{4.0}\def\xC{8.4}\def\xD{12.2}
  \node[actor] at (\xA,0) {Amministratore};
  \node[actor] at (\xB,0) {Orchestrator};
  \node[actor,fill=green!8] at (\xC,0) {Controller QKD};
  \node[actor,fill=orange!10] at (\xD,0) {Controller OTN};
  \foreach \x in {\xA,\xB,\xC,\xD} \draw[gray] (\x,-0.4) -- (\x,-9.0);
  \draw[msg] (\xA,-1.0) -- node[above]{richiesta servizio (SAE, QoS)} (\xB,-1.0);
  \draw[msg] (\xB,-1.5) -- ++(0.4,0) -- ++(0,-0.4) -- (\xB,-1.9);
  \node[right,align=left] at (\xB+0.45,-1.7) {address matching\\nodo OTN $\leftrightarrow$ nodo QKD};
  \draw[blue!60,rounded corners] (\xB-0.3,-2.55) rectangle (\xC+0.3,-3.85);
  \node[blue!60,anchor=south west,font=\scriptsize\itshape] at (\xB-0.3,-2.55) {ciclo ripetibile};
  \draw[msg] (\xB,-3.0) -- node[above]{creazione service link + vincoli} (\xC,-3.0);
  \draw[ret] (\xC,-3.6) -- node[above]{\texttt{svc\_link\_id}, cammini candidati} (\xB,-3.6);
  \draw[msg] (\xB,-4.6) -- node[above]{deploy del cammino scelto (\texttt{path\_id})} (\xC,-4.6);
  \draw[msg] (\xC,-5.0) -- ++(0.4,0) -- ++(0,-0.4) -- (\xC,-5.4);
  \node[right,align=left] at (\xC+0.45,-5.2) {key relay lungo\\il cammino};
  \draw[ret] (\xC,-6.1) -- node[above]{chiavi pronte} (\xB,-6.1);
  \draw[msg] (\xB,-6.9) -- node[above]{provisioning applicazione (\texttt{qkd\_app})} (\xC,-6.9);
  \draw[ret] (\xC,-7.6) -- node[above]{\texttt{sdqkdn\_application\_deployed}} (\xB,-7.6);
  \draw[msg] (\xB,-8.6) -- node[above,fill=white,inner sep=1pt]{i nodi OTN possono prelevare le chiavi dai nodi QKD} (\xD,-8.6);
\end{tikzpicture}
\caption[Provisioning end-to-end di un servizio QKD tra rete QKD e OTN]{Provisioning end-to-end di un servizio QKD tra rete QKD e OTN
(rielaborazione semplificata di~\cite[Fig.~2]{qkd018}). L'interazione con
il controller OTN è fuori dallo scope della specifica.}
\label{fig:seq018}
\end{figure}


I passi sono:
\begin{enumerate}
  \item l'amministratore chiede all'orchestrator un servizio, indicando i
        SAE nei nodi OTN e la loro QoS;
  \item l'orchestrator individua i controller coinvolti ed effettua
        l'address matching; conosce già gli indirizzi dei nodi QKD grazie
        a topologia e inventario;
  \item chiede al controller QKD la creazione di un service link e il
        calcolo dei cammini, con i vincoli desiderati;
  \item il controller assegna un \texttt{svc\_link\_id} e restituisce i
        cammini candidati, con una priorità per ciascuno; il passo può
        ripetersi con vincoli diversi;
  \item l'orchestrator sceglie un cammino e ne chiede il deploy; i nodi
        generano e inoltrano chiavi lungo il cammino finché i due estremi
        sono pronti;
  \item ricevuta la conferma, l'orchestrator richiede il provisioning
        dell'applicazione sul service link e avvisa il controller OTN che i
        nodi OTN possono iniziare a ricevere chiavi dai nodi QKD indicati.
\end{enumerate}
Si noti il parallelo con il workflow del link virtuale della 015: lì il
controller calcolava e configurava il cammino da solo, qui lo
\emph{propone} e attende la scelta dell'orchestrator. Con più domini QKD
l'orchestrator ripete i passi 3–5 per ciascun dominio lungo il cammino
inter-dominio scelto, collegando i segmenti attraverso i nodi di bordo: è
ciò che fanno Sim et al.\ (§\ref{sec:sim}).

% =====================================================================
\subsection{Protocolli e sicurezza}
\label{sec:protocolli}
% =====================================================================

\paragraph{Trasporto (§8).}
Come la 015, lo standard separa il modello dei dati (YANG~1.1) dal
protocollo di trasporto, che può essere NETCONF~\cite{rfc6241} o
RESTCONF~\cite{rfc8040}. Con RESTCONF ogni risorsa è un URI, la lettura usa
\texttt{GET} e la configurazione \texttt{POST}, \texttt{PUT},
\texttt{PATCH} e \texttt{DELETE}, con codifica XML o JSON; le notifiche si
ricevono sottoscrivendo i relativi stream. Lo standard osserva che il
paradigma REST, senza stato, si presta bene a questo tipo di API. A titolo
di esempio, la lettura della topologia e l'invocazione di una RPC avrebbero
la forma:
\begin{lstlisting}
GET  /restconf/data/etsi-qkd-sdn-topology:qkd_network
POST /restconf/operations/etsi-qkd-sdn-connectivity:create-service-link-id-for-app
\end{lstlisting}

\paragraph{Sicurezza (§7).}
Il controller ha una vista centrale della rete e ne passa un sottoinsieme
all'orchestrator; le misure di protezione dell'interfaccia (autenticazione,
cifratura del canale, per esempio TLS sotto RESTCONF) sono lasciate
all'implementazione. Il principio fondamentale, ereditato dalla GS~QKD~015,
è che \textbf{le chiavi derivate da QKD non transitano mai né per il
controller né per l'orchestrator}: nei moduli YANG compaiono solo tassi di
chiave e l'indirizzo del KME. Un orchestrator compromesso può quindi
causare disservizi (cammini sbagliati, risorse esaurite, denial of service)
ma non leggere le chiavi. Lo standard raccomanda procedure accurate di
progettazione, installazione e personalizzazione per mitigare richieste
malevole e attacchi ai canali di controllo.

% =====================================================================
\section{L'analogia con ACTN (RFC 8453)}
\label{sec:actn}
% =====================================================================

La GS~QKD~018 non inventa un modello nuovo: applica alla QKD un pattern già
standard nelle reti di trasporto, il framework ACTN
(\emph{Abstraction and Control of TE Networks}) dell'IETF~\cite{rfc8453}.
ACTN prevede tre livelli: il \emph{Customer Network Controller} (CNC), il
\emph{Multi-Domain Service Coordinator} (MDSC) e i \emph{Provisioning
Network Controller} (PNC) di ciascun dominio, collegati dalle interfacce
CMI (CNC–MDSC) e MPI (MDSC–PNC). La corrispondenza con la GS~QKD~018 è
diretta (Tab.~\ref{tab:actn}).

\begin{table}[h]
\centering\small
\caption{Corrispondenza tra ACTN e la pila ETSI per la QKD.}
\label{tab:actn}
\begin{tabularx}{\textwidth}{@{}l l L@{}}
\toprule
\textbf{ACTN} & \textbf{ETSI QKD} & \textbf{Osservazione} \\
\midrule
CNC & amministratore / portale cliente & origina la richiesta di servizio \\
MDSC & SDN orchestrator & vista multi-dominio, decisione del cammino \\
PNC & SDN controller del dominio QKD & calcola i candidati, attiva i cammini \\
MPI & interfaccia GS~QKD~018 & topologia astratta, provisioning, notifiche \\
SBI & interfaccia GS~QKD~015 & controllo dei singoli nodi \\
TE topology astratta & \texttt{etsi-qkd-sdn-topology} & vista vendor-neutral \\
banda TE & SKR / ESKR & la risorsa da allocare è la chiave \\
\bottomrule
\end{tabularx}
\end{table}

La differenza sostanziale sta nella risorsa. In ACTN la banda di un link è
fissata dalla capacità del canale; in QKD l'SKR fluttua con le condizioni
fisiche e la chiave si accumula in buffer finiti, da cui viene consumata.
Per questo nella 018 notifiche come \texttt{\_perf\_update} e
\texttt{\_overloaded} e la politica di \texttt{path\_restoration} hanno un
ruolo più centrale che nelle reti di trasporto classiche.

% =====================================================================
\section{Casi d'uso reali}
\label{sec:casi}
% =====================================================================

\subsection{SK Telecom: orchestrazione multi-vendor su KOREN}
\label{sec:sim}

Sim et al.~\cite{sim2023} descrivono il primo dispiegamento commerciale
documentato di un orchestrator conforme alla GS~QKD~018, sul testbed della
rete di ricerca coreana KOREN (2020–2022). Due operatori QKD gestiscono due
reti con dispositivi e KMS di vendor diversi, ciascuna con il proprio
controller (QKDNC); un orchestrator (QKDNO) le coordina. I link usano BB84
con decoy state su tratte tra $39$ e $100$~km. Le due reti si toccano in
\emph{nodi di bordo} (per esempio Pangyo e Daejeon) che non contengono
moduli QKD ma due KMS, uno per operatore, interconnessi.

Quando due SAE appartengono a reti diverse, l'orchestrator calcola il
cammino end-to-end con un \emph{path computation engine}, scegliendo
\emph{a priori} il nodo di bordo da attraversare; ogni controller crea poi
il proprio segmento di service link. Per legare i segmenti gli autori
estendono il modello: introducono un tipo di applicazione
\emph{multi-segment} (\texttt{MS}), con un identificativo
\texttt{ms\_app\_id} comune ai domini e le informazioni
\texttt{server\_ms\_info} e \texttt{client\_ms\_info} sui KMS esterni, e una
RPC \texttt{app-registration-ms}. Il trasferimento delle chiavi al nodo di
bordo usa l'interfaccia KMS–KMS della GS~QKD~020~\cite{qkd020}. Il caso
mostra sia la validità dello schema 018 sia un suo limite: i tipi di
applicazione ereditati dalla 015 (\texttt{CLIENT} e \texttt{INTERNAL}) non
bastano per il multi-dominio, e servono estensioni. È un buon esempio
dell'utilità delle \texttt{identity} YANG discusse in
§\ref{sec:yang-anatomia}: un nuovo tipo si aggiunge senza riscrivere il
modello di partenza.

\subsection{Politecnico di Milano: controllo gerarchico con ONOS}

La tesi di Vicario~\cite{vicario2025} realizza un prototipo end-to-end di
controllo gerarchico per reti metropolitane: i controller di dominio sono
basati su ONOS e gestiscono i nodi tramite il modello GS~QKD~015, mentre
l'interfaccia verso l'alto è la GS~QKD~018 esposta via RESTCONF a un
orchestrator scritto in Python che aggrega la topologia di più domini. È
un'implementazione utile perché mostra l'intera pila
015~$\to$~018 in software aperto, con gli stessi moduli YANG descritti in
§\ref{sec:yang018}.

\subsection{MadQCI: la rete di Madrid}

MadQCI~\cite{madqci2024} è una rete QKD reale a Madrid con decine di moduli
di vendor diversi su oltre $130$~km, integrata con la rete ottica di
produzione. Il controllo usa uno stack SD-QKD basato sulla GS~QKD~015, che
si conferma lo standard più adottato. La GS~QKD~018 è citata come strumento
per l'orchestrazione inter-dominio e la coordinazione ottico-quantistica, ma
gli autori la descrivono come un'attività ancora in corso. È una
testimonianza onesta dello stato dell'arte: il controllo di dominio è
maturo, l'orchestrazione molto meno.

% =====================================================================
\section{Limiti e prospettive}
\label{sec:limiti}
% =====================================================================

\begin{itemize}
  \item \textbf{Adozione.} La 015 è usata in reti reali, la 018 quasi solo
        in testbed e prototipi~\cite{madqci2024}.
  \item \textbf{Multi-dominio incompleto.} Lo scenario normativo è
        QKD$+$OTN; il caso QKD$+$QKD richiede estensioni come quelle di
        Sim et al.\ e il coordinamento con la GS~QKD~020 per il passaggio
        delle chiavi tra KMS.
  \item \textbf{Metriche.} L'ESKR è un indicatore istantaneo e non tiene
        conto della chiave già accumulata nei buffer né della sua
        dinamica. Laschet et al.~\cite{laschet2026} propongono indicatori
        più ricchi per la selezione dei cammini.
  \item \textbf{Riservatezza della topologia.} Esporre la topologia
        interna all'orchestrator è accettabile dentro un operatore, meno
        tra operatori diversi; servono schemi di instradamento inter-dominio
        che preservino la riservatezza~\cite{laschet2026,interconnect2026}.
  \item \textbf{Sicurezza dell'interfaccia.} Lo standard non specifica i
        meccanismi di protezione del canale orchestrator–controller, che
        restano a carico dell'implementazione.
  \item \textbf{Algoritmi.} Il calcolo dei cammini non è standardizzato:
        è un vantaggio per l'innovazione~\cite{routing2024,mendez2026}, ma
        rende meno prevedibile il comportamento in scenari multi-vendor.
\end{itemize}

% =====================================================================
\section{Conclusioni}
\label{sec:conclusioni}
% =====================================================================

Letta in sequenza, la produzione ETSI sul controllo della QKD ha una
logica chiara. La GS~QKD~002 ha mostrato che la QKD utile è una rete, con
gestione delle chiavi, requisiti di servizio differenziati e integrazione
nella rete ottica dell'operatore. La GS~QKD~015 ha risposto con
l'astrazione del nodo SD-QKD: moduli come interfacce, un key manager unico,
link fisici e virtuali, applicazioni con QoS, un controller che registra,
pianifica e configura, tutto descritto in YANG. La GS~QKD~018 ha aggiunto
il livello di orchestrazione: ogni dominio diventa un grafo astratto, le
risorse di chiave vengono monitorate tramite l'inventario e il provisioning
segue uno schema in cui il controller propone e l'orchestrator decide. Il
disegno ricalca ACTN, riusa YANG, NETCONF e RESTCONF e mantiene il principio
di sicurezza essenziale: le chiavi non lasciano mai i nodi. Le esperienze di
KOREN, del Politecnico di Milano e di MadQCI mostrano che lo schema funziona
ma che l'orchestrazione QKD è ancora giovane. Il passo successivo è
estenderla davvero al multi-operatore, con metriche migliori e garanzie di
riservatezza tra domini.


\printbibliography[heading=bibintoc,title={Bibliografia}]

\end{document}
